%% 
%% Copyright 2007-2025 Elsevier Ltd
%% 
%% This file is part of the 'Elsarticle Bundle'.
%% ---------------------------------------------
%% 
%% It may be distributed under the conditions of the LaTeX Project Public
%% License, either version 1.3 of this license or (at your option) any
%% later version. The latest version of this license is in
%% http://www.latex-project.org/lppl.txt
%% and version 1.3 or later is part of all distributions of LaTeX
%% version 1999/12/01 or later.
%% 
%% The list of all files belonging to the 'Elsarticle Bundle' is
%% given in the file `manifest.txt'.
%% 
%% Template article for Elsevier's document class `elsarticle'
%% with numbered style bibliographic references
%% SP 2008/03/01
%% $Id: elsarticle-template-num.tex 272 2025-01-09 17:36:26Z rishi $
%%
\documentclass[preprint,12pt]{elsarticle}

%% Use the option review to obtain double line spacing
%% \documentclass[authoryear,preprint,review,12pt]{elsarticle}

%% Use the options 1p,twocolumn; 3p; 3p,twocolumn; 5p; or 5p,twocolumn
%% for a journal layout:
%% \documentclass[final,1p,times]{elsarticle}
%% \documentclass[final,1p,times,twocolumn]{elsarticle}
%% \documentclass[final,3p,times]{elsarticle}
%% \documentclass[final,3p,times,twocolumn]{elsarticle}
%% \documentclass[final,5p,times]{elsarticle}
%% \documentclass[final,5p,times,twocolumn]{elsarticle}

%% For including figures, graphicx.sty has been loaded in
%% elsarticle.cls. If you prefer to use the old commands
%% please give \usepackage{epsfig}

%% The amssymb package provides various useful mathematical symbols
\usepackage{amssymb}
%% The amsmath package provides various useful equation environments.
\usepackage{amsmath}
%% The amsthm package provides extended theorem environments
%% \usepackage{amsthm}
\usepackage{xcolor}
%% The lineno packages adds line numbers. Start line numbering with
%% \begin{linenumbers}, end it with \end{linenumbers}. Or switch it on
%% for the whole article with \linenumbers.
\usepackage{lineno}
\usepackage{subcaption}
\usepackage{hyperref}
\usepackage{cleveref}
\usepackage{booktabs}
\usepackage{tabularx}
\usepackage{array}
\usepackage{siunitx}

\journal{Spectrochimica Acta Part A: Molecular and Biomolecular Spectroscopy}

\begin{document}

\begin{frontmatter}

%% Title, authors and addresses

%% use the tnoteref command within \title for footnotes;
%% use the tnotetext command for theassociated footnote;
%% use the fnref command within \author or \affiliation for footnotes;
%% use the fntext command for theassociated footnote;
%% use the corref command within \author for corresponding author footnotes;
%% use the cortext command for theassociated footnote;
%% use the ead command for the email address,
%% and the form \ead[url] for the home page:
%% \title{Title\tnoteref{label1}}
%% \tnotetext[label1]{}
%% \author{Name\corref{cor1}\fnref{label2}}
%% \ead{email address}
%% \ead[url]{home page}
%% \fntext[label2]{}
%% \cortext[cor1]{}
%% \affiliation{organization={},
%%    addressline={},
%%    city={},
%%    postcode={},
%%    state={},
%%    country={}}
%% \fntext[label3]{}

\title{Label-free Raman spectroscopic profiling of circulating extracellular
 vesicles in ischemic stroke: subject-level compositional mapping and the
 limits of individual clinical prediction}

%% use optional labels to link authors explicitly to addresses:
%% \author[label1,label2]{}
%% \affiliation[label1]{organization={},
%%    addressline={},
%%    city={},
%%    postcode={},
%%    state={},
%%    country={}}
%%
%% \affiliation[label2]{organization={},
%%    addressline={},
%%    city={},
%%    postcode={},
%%    state={},
%%    country={}}

\author[1]{[Author 1]\corref{cor1}} %% Author name
\ead{[author1@email.com]}
\cortext[cor1]{Corresponding author}

\author[1,2]{[Author 2]}

%% Author affiliation
\affiliation[1]{organization={[Department, Institution]},
   city={[City]},
   country={[Country]}}
   
\affiliation[2]{organization={[Department, Institution]},
   city={[City]},
   country={[Country]}}

%% Abstract
\begin{abstract}
Extracellular vesicles (EVs) reflect systemic physiological states and hold promise as liquid biopsy analytes. Raman spectroscopy provides a powerful, label-free method for profiling the complex biomolecular cargo of these EVs. However, translating Raman-based EV profiling to clinical settings is frequently hindered by statistical pseudoreplication, which treats multiple pixel-level spectra from the same individual as independent, thereby artificially inflating statistical power. To address this, we developed a subject-level biological principal component analysis (PCA) pipeline and applied it to Raman spectra (fingerprint and high-wavenumber C\textendash H stretching regions) of circulating EVs from 97 ischemic stroke patients and 50 healthy controls. Data were rigorously preprocessed and hierarchically aggregated to yield a single representative spectrum per participant. Subject-level PCA revealed that inter-individual spectral variance is overwhelmingly dominated by lipid compositional differences, explaining up to 80.4\% of the variance. Furthermore, intra-disease correlation analyses and cross-validated multivariate regression established that these bulk EV Raman signatures are not predictive of individual cognitive trajectories or clinical endpoints in the stroke cohort. We conclude that bulk circulating EV Raman profiles primarily capture underlying lipid variability, highlighting the necessity for affinity-based target enrichment or single-vesicle analysis in future translational applications.
\end{abstract}

%%Graphical abstract
% \begin{graphicalabstract}
% \includegraphics{grabs}
% \end{graphicalabstract}

% %%Research highlights
% \begin{highlights}
% \item Research highlight 1
% \item Research highlight 2
% \end{highlights}

%% Keywords
\begin{keyword}
Raman spectroscopy \sep extracellular vesicles \sep stroke \sep principal component analysis

%% keywords here, in the form: keyword \sep keyword

%% PACS codes here, in the form: \PACS code \sep code

%% MSC codes here, in the form: \MSC code \sep code
%% or \MSC[2008] code \sep code (2000 is the default)

\end{keyword}

\end{frontmatter}

%% Add \usepackage{lineno} before \begin{document} and uncomment 
%% following line to enable line numbers
%% \linenumbers

%% main text
%%

\section{Introduction}

The diagnostic utility of extracellular vesicles (EVs) is widely recognized, as their molecular cargo reflects disease-associated alterations. Raman spectroscopy offers a rapid, label-free approach for probing EV composition and detecting disease-related changes in their biochemical profiles. Recent work has reviewed the use of Raman spectroscopy for EV analysis \cite{liu2025} and highlighted the value of 2D and 3D Raman imaging for studying EVs \cite{horgan2020}. Methodological advances increasingly leverage data-driven approaches, including self-supervised learning to identify EVs from their Raman spectra \cite{jensen2024} and convolutional neural networks to classify EV spectra for prostate cancer detection \cite{lee2020}. Enhanced sensitivity has also been achieved using surface-enhanced Raman scattering (SERS), enabling EV-based cancer diagnostics \cite{koster2021}. Beyond oncology, Raman spectroscopy has been applied to detect amyloid beta in small EVs \cite{imanbekova2021} and to stratify patients with chronic kidney disease in type 2 diabetes using urinary EVs \cite{kaminska2022}. Raman-based characterization of EVs from bovine placenta and peripheral blood mononuclear cells has been further explored to develop diagnostic and therapeutic tools for the diagnosis and treatment of pregnancy complications \cite{zhang2020}.

\section{Related Work}

\color{teal}
\section{Materials and Methods}
\label{sec:methods}

\subsection{Study cohort and ethics}
\label{subsec:cohort}

The study cohort included 97 patients with ischemic stroke (45 females, 52 males) and 50 healthy controls (40 females, 10 males). Clinical cognitive function was assessed using the Montreal Cognitive Assessment (MoCA) at hospital admission (baseline) and at a 3-month follow-up. The longitudinal change in cognitive performance was calculated as $\Delta\text{MoCA} = \text{MoCA}_{\text{3m}} - \text{MoCA}_{\text{baseline}}$. Additional neuropsychological and radiological evaluations comprised the Repeatable Battery for the Assessment of Neuropsychological Status (RBANS) and the Alberta Stroke Program Early Computed Tomography Score (ASPECTS). The study protocol received approval from the relevant Institutional Review Board, and all participants or their legal representatives provided written informed consent before enrollment.

\subsection{Extracellular vesicle isolation}
\label{subsec:ev_isolation}

Circulating EVs were isolated from [plasma/serum] samples. After sequential centrifugation to remove cells and cellular debris, the resulting supernatant was subjected to size-exclusion chromatography (SEC). Fractions enriched in EVs were pooled based on protein concentration and particle count profiles, as determined by nanoparticle tracking analysis (NTA). The isolated EVs were stored at $-80\,^{\circ}\mathrm{C}$ until spectroscopic analysis.

\subsection{Raman spectroscopy data acquisition}
\label{subsec:acquisition}

Raman measurements were performed by depositing EV samples onto [CaF$_2$ windows] and allowing them to dry under controlled conditions. Spectra were obtained using a [Instrument Model] equipped with a [Wavelength] nm excitation laser and a [Magnification]$\times$ objective lens. For each participant, two spectral blocks were collected: the fingerprint region (770–1800~\si{\per\centi\metre}) and the high-wavenumber C\textendash H stretching region (xxx~\si{\per\centi\metre}). Within each block, spectral maps consisting of multiple pixels were acquired from distinct physical areas on each sample slide. Data were exported as Renishaw\texttrademark{} text files and imported into a structured binary data store (compressed NumPy arrays) prior to preprocessing.

\subsection{Data preprocessing pipeline}
\label{subsec:preprocessing}

To ensure robustness and eliminate artefactual variations prior to statistical analysis, a rigorous, multi-stage data preprocessing pipeline was developed. This pipeline operates sequentially through the four-level physical hierarchy of the measurement, progressing from the subject level to the sample slide, the physical area (spot), and finally the individual pixels. All preprocessing steps were implemented in Python and applied identically across the raw spectral hypercubes.

The procedure began with a quality control screen to filter out non-informative or corrupted spectra. Dead pixels exhibiting minimal signal variation or detector saturation were excluded. Cosmic-ray artifacts were identified and suppressed by comparing each pixel with its local spatial neighborhood and replacing extreme outliers with the local median. Spectra were required to meet a minimum signal-to-noise ratio within defined anchor windows to ensure sufficient spectral quality. The boundaries of both spectral blocks were trimmed to eliminate edge-related noise from the spectrometer. Areas with inhomogeneous sample deposition, characterized by pronounced coffee-ring effects, were also removed.

After quality control, the remaining spectra underwent baseline correction to remove the broad fluorescence background without distorting underlying Raman bands. The asymmetric reweighted penalized least-squares (arPLS) algorithm was used for this purpose, followed by a secondary spike-repair pass to address any residual sharp anomalies. To improve signal quality, spectra were smoothed using region-specific Savitzky-Golay filters and interpolated onto a uniform wavenumber grid. In the final step, each spectral block was normalized independently to account for intensity variations. The fingerprint region was transformed using Standard Normal Variate (SNV) to correct multiplicative scatter, while the C\textendash H stretch region was area-normalized to preserve the relative proportions of lipid and protein band intensities. A summary of the specific parameters used at each stage of the preprocessing pipeline is provided in \cref{tab:preprocessing_params}.

\begin{table}[htbp]
\centering
\caption{Summary of parameters used in the multi-stage Raman spectral preprocessing pipeline.}
\label{tab:preprocessing_params}

\begin{tabularx}{\textwidth}{@{}>{\raggedright\arraybackslash}p{0.34\textwidth}
>{\raggedright\arraybackslash}X@{}}
\toprule
\textbf{Preprocessing Stage} & \textbf{Parameters} \\
\midrule

\textit{Quality Control} & \\
Dead-pixel rejection & Minimum SD: 1.0; Maximum intensity: 65,000 \\
Cosmic-ray suppression & $z$-score threshold: 8.0; Spatial window: radius 2 pixels \\
SNR filtering & Minimum SNR: 5.0, relative to flat-segment noise \\
Edge trimming & 10 channels, $\approx$15~\si{\per\centi\metre}, from each block boundary \\
Coffee-ring screen & Intensity-ratio threshold: 2.0 \\
Minimum survival & $\ge$ 4 pixels per spot; $\ge$ 10 pixels per subject \\

\midrule
\textit{Baseline Correction} & \\
arPLS algorithm & $\lambda = 10^{6}$, $\rho = 10^{-3}$, maximum 100 iterations \\
Spike repair & Whitaker--Hayes $z$-score: 8.0; window: 5 channels; max 2\% repair \\

\midrule
\textit{Smoothing \& Interpolation} & \\
Savitzky--Golay (Fingerprint) & Window: 7 channels; polynomial order: 2 \\
Savitzky--Golay (C\textendash H stretch) & Window: 9 channels; polynomial order: 2 \\
Interpolation & Cubic spline; step size: 1.0~\si{\per\centi\metre} \\
\midrule
\textit{Normalisation} & \\
Fingerprint region & Standard Normal Variate (SNV) \\
C\textendash H stretch region & Area integration to unity, trapezoidal rule \\
\bottomrule
\end{tabularx}
\end{table}


\subsection{Subject-level biological PCA}
\label{subsec:biopca}

To address pseudoreplication, preprocessed pixel-level spectra were hierarchically aggregated within each spectral block before conducting cross-subject analyses. For each subject, the median value was calculated across all surviving pixels in the sample, regardless of their physical origin, resulting in a single representative spectrum per subject per spectral block. Median aggregation was selected instead of the mean due to its greater robustness against residual outlier pixels that may persist after quality filtering.

Three independent PCA models were fitted to the resulting subject-level median spectra matrices: \textsc{PCA\_{1330}} (fingerprint block, $n \times 986$), \textsc{PCA\_{2850}} (C\textendash H stretch block, $n \times 975$), and \textsc{PCA\_{whole}} (unweighted column-wise concatenation of both blocks, $n \times 1961$). PCA was performed using the original normalized intensity scale, without additional variance scaling. For each model, the first three principal components (PC1, PC2, PC3) were extracted. Dominant loading regions for each component were identified by applying a 90th-percentile absolute-loading threshold, merging contiguous intervals separated by less than 6~\si{\per\centi\metre}, and retaining intervals wider than 8~\si{\per\centi\metre}. Biochemical assignments were based on published Raman band tables for lipids, proteins, and nucleic acids~\cite{kaminska2022,imanbekova2021}.

\subsection{Intra-disease clinical correlation analysis}
\label{subsec:correlation}

Within the stroke cohort, Spearman's rank correlation coefficient ($\rho$) was calculated between subject-level normalized intensity at each wavenumber position and five clinical endpoint variables: MoCA baseline score, MoCA 3-month score, $\Delta$MoCA, RBANS total score, and ASPECTS. Age was also included as a correlate. Pearson's $r$ was used instead of Spearman's $\rho$ when the Shapiro--Wilk normality test ($\alpha = 0.05$) indicated a Gaussian distribution for the endpoint variable.

To control the family-wise error rate, cluster-mass permutation testing was performed using 1000 permutations and a cluster-defining threshold of $\alpha = 0.05$. Clusters were defined as contiguous runs of adjacent wavenumbers, each exceeding the per-comparison threshold. The test statistic was the sum of the absolute Spearman coefficients within each cluster (cluster mass). The empirical null distribution of the maximum cluster mass was generated by permuting the endpoint vector 1000 times while keeping the spectral matrix fixed. All tests were two-tailed. The Benjamini--Hochberg FDR correction was applied separately across all wavenumbers within each endpoint--block combination ($\alpha_{\text{FDR}} = 0.05$). A spectral interval was considered significant only if it passed both cluster-mass permutation testing and FDR correction.

\subsection{Targeted biochemical band analysis}
\label{subsec:band_analysis}

Thirteen Raman features were defined \textit{a priori} based on established biochemical assignments, spanning the fingerprint and high-wavenumber regions. For each feature, the Raman signal was summarized as the trapezoidal integral within a predefined spectral window, thereby avoiding data-driven peak selection. In addition to these absolute band integrals, six cross-region ratios were computed to capture relative contributions from protein-, lipid-, nucleic-acid-, and CH-stretch-associated signals. These ratios were amide I/symmetric CH$_2$ stretch, amide I/symmetric CH$_3$ stretch, phenylalanine/symmetric CH$_3$ stretch, DNA phosphate/symmetric CH$_2$ stretch, amide III/amide I, and symmetric CH$_2$/symmetric CH$_3$ stretch. Spearman's $\rho$ was then calculated between each band integral or ratio and the three MoCA endpoints. To control for multiple testing across the full feature set, $p$-values were adjusted using the Benjamini--Hochberg false-discovery-rate procedure with $\alpha_{\mathrm{FDR}} = 0.05$.



\begin{table}[ht]
\centering
\caption{Predefined Raman bands used for integral-based feature extraction.}
\label{tab:raman_bands}
\begin{tabular}{ll}
\hline
Biochemical assignment & Spectral window / \si{\per\centi\metre} \\
\hline
Phenylalanine ring stretch & 995--1015 \\
DNA phosphate symmetric stretch & 1080--1100 \\
C--C skeletal stretch & 1115--1140 \\
Amide III & 1240--1280 \\
CH$_2$ deformation & 1290--1320 \\
CH$_2$ scissoring & 1435--1470 \\
Carboxylate stretch & 1565--1595 \\
Amide I & 1630--1680 \\
Symmetric CH$_2$ stretch & 2840--2865 \\
Asymmetric CH$_2$ stretch & 2875--2900 \\
Symmetric CH$_3$ stretch & 2920--2945 \\
Asymmetric CH$_3$ stretch & 2950--2975 \\
O--H/N--H stretching & 3045--3080 \\
\hline
\end{tabular}
\end{table}

\subsection{Multivariate predictive modelling}
\label{subsec:pls}

Cross-validated partial least squares (PLS) regression was conducted on the concatenated subject-level spectral matrix, comprising fingerprint and C–H stretch blocks, to determine whether a latent multivariate spectral signature predicts MoCA scores. The number of PLS components was optimized using inner cross-validation across the grid $\{1, 2, 3, 4, 5\}$, while the outer loop implemented 5-fold KFold cross-validation. Model performance was evaluated using cross-validated $R^{2}$ and root mean squared error (RMSE). Statistical significance was determined by permuting the endpoint vector 1000 times; the empirical $p$-value was defined as the proportion of permuted $R^{2}$ values exceeding the observed value. A PLS permutation significance threshold ($p_{\text{perm}} < 0.05$) was required before fitting any sparse PLS discriminant analysis (sPLS-DA) classifier for binarised MoCA endpoints.

A secondary minimal candidate panel was developed using nested cross-validated ElasticNet regression (5 outer by 3 inner folds, $n_{\text{bootstrap}} = 200$, $n_{\text{permutations}} = 1000$), targeting the selection of 5 to 10 spectral features per endpoint.

\color{black}


\section{Results and Discussion}
\subsection{Spectral characteristics of circulating EVs and severity-associated variations}
\label{subsec:spectral_profiles}

Figure~\ref{fig:sp_mean} shows the mean Raman spectra of EVs from the disease and control cohorts, and the corresponding band assignments are summarized in Table~\ref{tab:assignment}. The spectra exhibit the classical vibrational profile of a complex, lipid- and protein-rich biological matrix, consistent with EV composition and featuring bands attributable to nucleic acids, proteins, lipids, and carotenoids.

For nucleic acids, characteristic bands include the cytosine/uracil breathing and stretching modes at 786~cm$^{-1}$, the O--P--O backbone stretch at 821~cm$^{-1}$, and a sugar-related C--O--C stretch at 925~cm$^{-1}$. In addition, the band at 1112~cm$^{-1}$ reflects PO$_2^-$ stretching contributions that can overlap with protein-related C--N stretching.

For proteins, the observed bands include the C--C stretch associated with $\beta$-sheet structure at 980~cm$^{-1}$, the C--O/C--N stretching region at 1060~cm$^{-1}$, and the amide III band peaking at 1258~cm$^{-1}$ (which overlaps with CH bending in lipids). Aromatic amino-acid features are also included, such as the tyrosine doublet at 851~cm$^{-1}$ (overlapping with glycan contributions), tryptophan at 1555~cm$^{-1}$, and the amide I band at 1657~cm$^{-1}$ (overlapping with lipid C=C stretching). Protein–lipid overlap is also represented in the high-wavenumber CH stretching bands at 2890 and 2937~cm$^{-1}$.

Lipid signatures include C--C stretching around 1070--1100cm$^{-1}$, CH bending/amide III overlap at 1200--1270cm~$^{-1}$, CH$_2$ twisting at 1301~cm$^{-1}$, and CH$_2$/CH$_3$ scissoring at 1440~cm$^{-1}$, and CH$_2$/CH$_3$ bending at 1465~cm$^{-1}$. Carbonyl stretches are represented at 1721~cm$^{-1}$ and 1744~cm$^{-1}$ (tentatively assigned to more and less hydrogen-bonded C=O groups, respectively). Finally, lipid-rich CH stretching bands are included at 2853~cm$^{-1}$ (CH$_2$ symmetric stretch) and, together with protein overlap, at 2890 and 2937~cm$^{-1}$.

For carotenoids, two main bands were observed: the C--C stretch at 1156~cm$^{-1}$ and the conjugated --C=C-- stretch at 1525~cm$^{-1}$.

Notably, most peaks showed little frequency shift between cohorts; however, their relative intensities differed markedly. The high intensities of the bands at 1465 and 2850~cm$^{-1}$ indicate that the EVs under study are lipid-rich.

The most pronounced cohort differences were observed in lipid-associated bands. In particular, the bands at 1258, 1301, 1440, 1657, and 2853~cm$^{-1}$ exhibited lower relative intensities in the disease cohort. The 1430/1465~cm$^{-1}$ intensity ratio also differed between the control and disease cohorts, which may indicate differences in fatty-acid composition \citenum{czamara2015}. Disease-related alterations in EV lipid composition have been reported previously \cite{ghadami2023, fyfe2023, zhang2025}; however, this observation should be interpreted with caution, as these lipid peaks varied markedly across individuals (Fig.~Sx).

Both carotenoid-associated bands were more intense in the control cohort. Consistent with this, Ref.~\citenum{gualerzi2025} reported an increase in EV carotenoid content following rehabilitation in Parkinson's disease, potentially linked to patient redox status.


%\begin{figure}[htbp]
%\centering
%\includegraphics[width=\linewidth]{figures/fig_mean_spectra_1300.png}
%\caption{Group mean Raman spectra in the fingerprint region (770--1800~\si{\per\centi\metre}). Solid lines represent the group mean; shaded envelopes indicate one standard deviation across all subjects. Key band assignments are annotated.}
%\label{fig:sp_mean_1300}
%\end{figure}

%\begin{figure}[htbp]
%\centering
%\includegraphics[width=\linewidth]{figures/fig_mean_spectra_2850.png}
%\caption{Group mean Raman spectra in the high-wavenumber C\textendash H stretching region (2400--3200~\si{\per\centi\metre}). Solid lines represent the group mean; shaded envelopes indicate one standard deviation.}
%\label{fig:sp_mean_2850}
%\end{figure}


\begin{figure}[htbp]
\centering
\includegraphics[width=\linewidth]{figures/mean_spectra.eps}
\caption{Mean Raman spectra for the control (black) and disease (purple) groups. Colored regions highlight Raman bands predominantly associated with lipids, proteins, nucleic acids, or carotenoids; gray regions indicate substantial contributions from two biomolecular classes. Band assignments are summarized in Table~\ref{tab:assignment}. Spectral standard deviations are shown in Fig.~Sx.}
\label{fig:sp_mean}
\end{figure}


\begin{table}[htbp]
\centering
\scriptsize
\caption{Raman band assignments for EV spectra.}
\label{tab:assignment}
\begin{tabularx}{\textwidth}{XXXXXX}
\hline
Vibration & Biochemical assignment & Peak / \si{\per\centi\metre} & Range / \si{\per\centi\metre} & Reference \\
\hline
Cytosine/uracil breathing/stretching modes & nucleic acids & 786 & 773--796 & \citenum{hobro2007} \\
O--P--O stretch & nucleic acids & 821 & 812--832 & \citenum{buchan2025, ferrara2015}  \\
Tyrosine doublet/glycan signature & proteins/glycans & 851 & 836--874 & \citenum{gualerzi2017, grenier2025, buchan2025} \\
C--O--C stretch in sugars & nucleic acids & 925 & 900--938 & \citenum{zhou2026}  \\
C--C stretch $\beta$-sheet & proteins & 980 & 965--995 & \citenum{grenier2025} \\
C--O/C--N stretch & proteins & 1060 & 1044--1078 & \citenum{gualerzi2017} \\
C--C stretch & lipids & 1089 & 1082--1098 & \citenum{potcoava2021} \\
PO$_2^-$ stretch/C--N stretch & nucleic acids/proteins & 1112 & 1105--1120 & \citenum{gualerzi2019, grenier2025} \\
C--C stretch carotene & carotenoids & 1156 & 1145--1165 & \citenum{grenier2025} \\
CH bend/amide III & lipids/proteins & 1258 & 1240--1280 & \citenum{potcoava2021, grenier2025} \\
CH$_2$ twist & lipids & 1301 & 1290--1325 & \citenum{grenier2025} \\
CH$_2$/CH$_3$ scissor deformation & lipids & 1440 & 1430--1445 & \citenum{czamara2015} \\
CH$_2$/CH$_3$ bend & lipids & 1465 & 1450--1475 & \citenum{czamara2015} \\
--C=C-- stretch carotene & carotenoids & 1525 & 1510--1535 & \citenum{gualerzi2017, grenier2025} \\
Tryptophan signature & proteins & 1555 & 1550--1560 & \citenum{gualerzi2017} \\
C=C stretch/amide I & lipids/proteins & 1657 & 1640--1675 & \citenum{gualerzi2019, potcoava2021} \\
C=O stretch (hydrogen bonded) & lipids & 1721 & 1710--1730 & \citenum{dreier2019} \\
C=O stretch & lipids & 1744 & 1735--1755 & \citenum{grenier2025, potcoava2021} \\
CH$_2$ symmetric stretch & lipids & 2853 & 2845--2860 & \citenum{czamara2015, gualerzi2017} \\
CH$_2$ asymmetric stretch & lipids/proteins & 2890 & 2870--2910 & \citenum{czamara2015, gualerzi2017} \\
CH$_3$ asymmetric \& asymmetric stretch (overlapping) & lipids/proteins & 2937 & 2925--2945 & \citenum{czamara2015, gualerzi2017} \\
\hline
\end{tabularx}
\end{table}




%\cref{fig:sp_mean_1300} and \cref{fig:sp_mean_2850} display the mean Raman spectra of circulating EVs in the fingerprint and high-wavenumber C\textendash H stretching regions, respectively. Both spectral domains exhibit the classical vibrational profile of a complex, lipid- and protein-rich biological matrix. In the fingerprint region (\cref{fig:sp_mean_1300}), the principal spectral features include the phenylalanine ring breathing mode at $\approx$1003~\si{\per\centi\metre}, the DNA/RNA phosphodiester symmetric stretch near $\approx$1090~\si{\per\centi\metre}, the C\textendash C skeletal stretch at $\approx$1127~\si{\per\centi\metre}, the amide~III envelope (1200--1270~\si{\per\centi\metre}), the CH$_2$ twisting deformation ($\approx$1300~\si{\per\centi\metre}), the CH$_2$/CH$_3$ scissoring/deformation complex ($\approx$1450~\si{\per\centi\metre}), and the amide~I band ($\approx$1657~\si{\per\centi\metre})~\cite{kaminska2022}. In the C\textendash H stretching region (\cref{fig:sp_mean_2850}), the dominant features arise from the symmetric and asymmetric CH$_2$ and CH$_3$ stretching vibrations (2850--2940~\si{\per\centi\metre}) originating from lipid acyl chains and protein methyl groups.

%\begin{figure}[htbp]
%\centering
%\includegraphics[width=\linewidth]{figures/fig_mean_spectra_1300.png}
%\caption{Group mean Raman spectra in the fingerprint region (770--1800~\si{\per\centi\metre}). Solid lines represent the group mean; shaded envelopes indicate one standard deviation across all subjects. Key band assignments are annotated.}
%\label{fig:sp_mean_1300}
%\end{figure}

%\begin{figure}[htbp]
%\centering
%\includegraphics[width=\linewidth]{figures/fig_mean_spectra_2850.png}
%\caption{Group mean Raman spectra in the high-wavenumber C\textendash H stretching region (2400--3200~\si{\per\centi\metre}). Solid lines represent the group mean; shaded envelopes indicate one standard deviation.}
%\label{fig:sp_mean_2850}
%\end{figure}

To determine whether ischemic stroke severity imprints a distinct spectral signature upon circulating EVs, we computed the difference spectra between the severe and mild stroke cohorts (\cref{fig:diff_spectrum}). While the normalized mean intensities of the two groups overlap extensively upon visual inspection, an analysis of effect sizes (Cohen's $d$) reveals specific compositional shifts driven by stroke severity. Within the fingerprint region, EVs isolated from patients with severe stroke exhibit relative intensity enhancements in the amide~III region (1200--1300~\si{\per\centi\metre}) and the CH$_2$ scissoring/deformation band ($\approx$1450~\si{\per\centi\metre}). These spectral deviations indicate structural alterations in protein secondary structure and lipid acyl chain packing. Conversely, the mild stroke cohort demonstrates higher relative intensities in the C\textendash C skeletal stretch region (1050--1150~\si{\per\centi\metre}) and the 1500--1600~\si{\per\centi\metre} domain.

In the high-wavenumber region, the spectral deviations are more subtle; however, severe stroke is associated with positive effect sizes proximate to the primary symmetric and asymmetric CH$_2$ stretching modes (2850--2900~\si{\per\centi\metre}). These severity-dependent variations suggest that, although the global spectral architecture remains highly conserved, pronounced clinical severity induces specific alterations in the lipid and protein organization of the circulating EV pool.

\begin{figure}[htbp]
\centering
\includegraphics[width=\linewidth]{figures/fig_difference_spectrum.png}
\caption{Difference spectrum analysis comparing the severe and mild stroke groups. The upper panels depict the normalized mean Raman intensity for the fingerprint (left) and C\textendash H stretching (right) regions. The lower panels display the signed effect size ($\text{sign}(\Delta) \times |\text{Cohen's } d|$) across the spectral range, delineating specific wavenumbers associated with disease severity.}
\label{fig:diff_spectrum}
\end{figure}

\subsection{Spectral correlations with clinical endpoints and cognitive trajectories}
\label{subsec:correlations}

To further interrogate the relationship between EV biomolecular cargo and clinical outcomes, we mapped the correlation landscape across the spectral features. \cref{fig:corr_heatmap} presents the Spearman correlation heatmaps for the fingerprint and C\textendash H stretching regions against key clinical and demographic variables.

Rather than exhibiting a uniform global alteration, the heatmaps reveal a highly localized correlational architecture. For the cognitive assessments, the baseline and 3-month Montreal Cognitive Assessment (MoCA) scores demonstrate relatively diffuse correlation patterns. Crucially, however, the longitudinal trajectory of cognitive change ($\Delta$MoCA) yields a distinct cluster of positive correlation in the fingerprint region, centered squarely on the amide~III band (1214--1364~\si{\per\centi\metre}). This observation parallels the severity-associated increases identified in the difference spectra, implying that protein conformational changes within the EV pool may co-vary with the biological mechanisms driving cognitive decline.

Furthermore, the Repeatable Battery for the Assessment of Neuropsychological Status (RBANS) total score exhibits discrete clusters of correlation. Positive associations emerge near 1510~\si{\per\centi\metre}---potentially overlapping with carotenoid C=C stretching---and within the 2895--2998~\si{\per\centi\metre} domain in the high-wavenumber region, which captures the CH$_3$ stretching commonly associated with protein methyl groups.

Importantly, baseline demographics---particularly age---profoundly sculpt the spectral landscape. Age exhibits pronounced correlation clusters, notably a robust positive band extending from 1510 to 1651~\si{\per\centi\metre}, interspersed with alternating negative regions. This pervasive influence underscores that systemic metabolic and demographic factors are major drivers of EV compositional heterogeneity.

\begin{figure}[htbp]
\centering
\includegraphics[width=\linewidth]{figures/fig_clinical_corr_heatmap.png}
\caption{Spearman correlation heatmaps delineating the relationship between subject-level Raman spectral features and clinical endpoints. The upper panel illustrates the fingerprint region (733--1789~\si{\per\centi\metre}), while the lower panel displays the C\textendash H stretching region (2450--3198~\si{\per\centi\metre}). Color intensity corresponds to the magnitude and direction of the Spearman rank correlation coefficient ($\rho$).}
\label{fig:corr_heatmap}
\end{figure}

\subsection{Biochemical context and implications for clinical translation}
\label{subsec:biochem}

Direct inspection of the spectral data establishes that the circulating EV pool harbors severity-associated and cognition-associated biochemical signatures, most notably within the amide~III, CH$_2$/CH$_3$ stretching, and C=C regions. Nevertheless, the overwhelming variability observed in bulk EV preparations is driven by systemic lipid heterogeneity. This conclusion is strongly supported by the dominant CH$_2$ features and the pervasive influence of patient age on the correlation profiles.

Although ischemic stroke precipitates the release of brain-derived EVs (BDEVs) carrying neuronal and glial cargo, these particles constitute only a minor fraction of the total circulating pool, which is heavily dominated by platelet-, erythrocyte-, and endothelial-derived vesicles~\cite{liu2025}. The severity-dependent intensity enhancements in structural bands ($\approx$1250, $\approx$1450, and 2850--2900~\si{\per\centi\metre}) and their correlation with $\Delta$MoCA likely reflect a generalized systemic inflammatory and metabolic response to cerebrovascular injury, rather than an isolated neuro-specific signal.

While these spectral correlations yield valuable biological insights, the relatively modest magnitude of these shifts underscores the fundamental limitations of employing bulk EV Raman spectroscopy for robust individual-level clinical prediction. Advanced multivariate predictive modeling previously applied to this dataset confirmed that these subtle features fail to overcome the systemic noise floor, yielding predictions that do not exceed cohort-mean estimates. Consequently, future translational efforts must focus on surmounting this biological dilution. Realizing reliable clinical predictions will unequivocally necessitate either the affinity-based enrichment of specific EV subpopulations (e.g., utilizing neural surface markers) prior to Raman acquisition~\cite{liu2025}, or the deployment of single-vesicle spectroscopic techniques~\cite{jensen2024, koster2021} to effectively isolate the target pathological signatures from the systemic background.

\section{Conclusions}
\label{sec:conclusions}

This study has presented a rigorous, subject-level Raman spectroscopic analysis of circulating EVs in 97 ischemic stroke patients and 50 healthy controls. The following conclusions are drawn from the data:

\begin{enumerate}
  \item A multi-stage preprocessing pipeline incorporating hierarchical median aggregation effectively eliminates pseudoreplication, enabling the robust characterization of inter-individual EV variability.
  \item Difference spectrum analysis reveals that severe ischemic stroke is associated with distinct compositional shifts in circulating EVs, notably an enhancement in the relative intensity of the amide~III region (1200--1300~\si{\per\centi\metre}) and the CH$_2$ scissoring/deformation bands ($\approx$1450~\si{\per\centi\metre}). This suggests alterations in protein secondary structure and lipid chain packing.
  \item Spectral correlation heatmaps demonstrate that the trajectory of cognitive change ($\Delta$MoCA) mirrors these severity-associated shifts, localizing primarily to the amide~III band. Age and other baseline demographics also strongly pattern the spectral landscape, underscoring the profound systemic metabolic influence on EV composition.
  \item Despite these targeted insights, the overwhelming inter-individual variance remains driven by systemic lipid heterogeneity. The severity-dependent signals likely reflect a generalized systemic inflammatory and metabolic response rather than an isolated neuro-specific signature.
  \item Future clinical applications of EV Raman spectroscopy in cerebrovascular disease must prioritize the affinity-based isolation of brain-derived EV subpopulations or the implementation of single-vesicle Raman analysis to overcome the profound signal dilution imposed by the dominant lipid background of non-cerebral circulating EVs.
\end{enumerate}


\section*{Data Availability}

The processed spectral data matrices and analysis scripts supporting the findings of this study are available from the corresponding author on reasonable request. Raw Raman data are subject to local data governance restrictions.

\section*{Ethics Declaration}

The study protocol received approval from the relevant Institutional Review Board ([IRB reference]). Written informed consent was obtained from all participants or their legal representatives.

\section*{Author Contributions}

[CRediT taxonomy statement to be completed by authors]

\section*{Declaration of Competing Interests}

The authors declare no competing interests.

\section*{Funding}

[Funding statement to be completed by authors]

\section*{Acknowledgements}

The authors thank [to be completed by authors].

\bibliographystyle{elsarticle-num} 
\bibliography{cas-refs}

\end{document}

\endinput
%%
%% End of file `elsarticle-template-num.tex'.
